\section{The simultaneous cubic-unit gate on a fixed elliptic curve}
\label{ssg-section}
Fix \(g=3,u=\zeta\). Write
\[
 A_0=s^3-3s-1,\quad B_0=3s(s+1),\quad A=-B_0,\quad B=A_0+B_0.
\]
The smooth projective normalization \(D=D_{3,\zeta}\) of
\begin{equation}\label{ssg-source}
 v^2=1+r(s),\qquad w^2=1-3r(s),\qquad r=A/B
\end{equation}
has genus six by HG. Its positive rational locus is
\(0<r<1/3,\ v>0\); its seed parameter is \(t=(v+w+2)/(v-w)>1\).
Both square roots must have the same \(s\). Their quotient coordinates are
\[
 y_1=Bv,\quad y_1^2=A_0(A_0+B_0),\qquad
 y_2=Bw,\quad y_2^2=(A_0+B_0)(A_0+4B_0).
\]

\begin{theorem}[A quartic over the fixed elliptic quotient]\label{ssg-quartic}
Put
\begin{align*}
 f&=X^3-9X-9,& h&=X+2,\\
 H&=f+18h^2=X^3+18X^2+63X+63,&
 R&=f-108h^2=X^3-108X^2-441X-441,\\
 J&=6Xh,& d&=4X+9.
\end{align*}
Over \(E:Y^2=f(X)\), the equation
\begin{equation}\label{ssg-equation}
 Z^4-2H(X)Z^2+f(X)R(X)=0
\end{equation}
has smooth projective normalization isomorphic over \(\mathbb Q\) to
\(D\). It is geometrically connected of genus six and has degree four over \(E\).
\end{theorem}
\begin{proof}
On \(s\ne-1\), define
\begin{align}
 X&=-\frac{2s^2+5s+2}{(s+1)^2},&
 Y&=\frac{y_1}{(s+1)^3},&
 Z&=\frac{y_2}{(s+1)^3},\label{ssg-forward}\\
 \kappa&=\frac{s-1}{s+1},&
 \gamma&=\frac{A_0+B_0}{(s+1)^3}.&&\notag
\end{align}
The first quotient identity gives \(Y^2=f\). Clearing denominators gives
\begin{align}
 \kappa^2&=d,& \frac{B_0}{(s+1)^3}&=-3h,
 &\gamma&=\frac{-X\kappa-3h}{2},\label{ssg-gamma}\\
 Y^2&=\gamma^2+3h\gamma,&
 Z^2&=\gamma^2-9h\gamma=H+J\kappa.&&\notag
\end{align}
Indeed \(A_0/(s+1)^3=\gamma+3h\). Multiplying the normalized
cubic factors proves the square formulas; their sum is \(-X\kappa\).
Expansion gives
\begin{equation}\label{ssg-discriminant}
 H^2-fR=J^2d.
\end{equation}
Eliminating \(\kappa\) proves the quartic.
Conversely, on the open with nonzero denominators, put
\begin{equation}\label{ssg-inverse}
 \kappa=\frac{Z^2-H}{J},\quad
 s=\frac{1+\kappa}{1-\kappa},\quad
 y_1=Y(s+1)^3,\quad y_2=Z(s+1)^3.
\end{equation}
The quartic and \eqref{ssg-discriminant} imply \(\kappa^2=d\).
The identities \eqref{ssg-gamma} recover the original two equations
at this same \(s\). The maps are inverse on nonempty dense opens.

The maps \(D\to H_1\) and \(H_1\to E\) both have degree two.
Thus \(\mathbb Q(D)/\mathbb Q(E)\) has degree four; the inverse
formulas show that \(Z\) generates it. The monic quartic is its
minimal polynomial. The same argument over the algebraic closure
proves geometric irreducibility. The function-field isomorphism
extends to the smooth projective models by \cite{StacksCurves2026},
Lemma 53.2.2 and Theorem 53.2.6.
\end{proof}

\begin{theorem}[The exact positive rational domain]\label{ssg-positive}
The positive rational locus of \(D\) is in bijection with the rational
triples satisfying
\begin{equation}\label{ssg-positive-test}
 Y^2=f(X),\quad Z^4-2HZ^2+fR=0,\quad
 0<Z^2<R(X),\quad Y(X+2)>0.
\end{equation}
This is a bijection of curve points with their power-map parameter,
not a claim of a unique such point above each integer seed at \(g=3\).
\end{theorem}
\begin{proof}
For a positive point, \(s=-1\) and \(s=\infty\) are impossible:
both have homogeneous \(r\)-value zero. Hence \eqref{ssg-forward}
is defined. Here \(r=3h/\gamma,\ v=Y/\gamma,\ w=Z/\gamma\), and
\begin{equation}\label{ssg-delta}
 \Delta:=Y^2-Z^2=12h\gamma,\qquad r=\frac{36h^2}{\Delta}.
\end{equation}
The value \(h=0\) forces \(s=0\) or \(\infty\) in the EQ
projective abscissa relation, again with \(r=0\). Thus \(0<r<1/3\)
implies \(\Delta>108h^2>0\) and \(Z^2<R\).
A rational positive point cannot have \(w=0\), since this would give
\(3v^2=4\). Therefore \(Z^2>0\). Since \(\gamma\) and \(h\) have
the same sign, positivity of \(v\) is exactly \(Yh>0\).

Conversely assume \eqref{ssg-positive-test}. The value \(X=0\) is
impossible because \(f(0)=-9\). At \(X=-2\), one has
\(f=H=R=1,J=0\); the quartic becomes \((Z^2-1)^2=0\), inconsistent
with \(0<Z^2<R\). Thus \(J\ne0\).
Formula \eqref{ssg-inverse} gives rational \(\kappa\) with
\(\kappa^2=4X+9\). Neither \(\kappa=1\) nor \(-1\) occurs,
since either would force \(X=-2\). The resulting \(s\) is finite
and not \(-1\), and both hyperelliptic equations hold there. Set
\begin{equation}\label{ssg-positive-inverse}
 \Delta=f-Z^2,\quad \gamma=\frac{\Delta}{12h},\quad
 v=\frac{Y}{\gamma},\quad w=\frac{Z}{\gamma}.
\end{equation}
The inequality gives \(\Delta>108h^2>0\).
Equations \eqref{ssg-gamma} recover the normalized \(B\)-coordinate
\(\gamma\), so \(r=36h^2/\Delta\in(0,1/3)\);
\(Yh>0\) makes \(v>0\). Equivalently the quartic rearranges to
\begin{equation}\label{ssg-conic}
 (f-Z^2)^2=36h^2(3f+Z^2),
\end{equation}
giving \(3v^2+w^2=4\). These constructions are inverse.

The point \(O\) of \(E\) corresponds to \(s=-1\), outside the positive
locus. The repeated-discriminant value \(X=-9/4\) is not in
\(E(\mathbb Q)\), since \(f(-9/4)=-9/64\).
Positivity avoids all HG branch values; normalization therefore
introduces no extra choice of local branch in this bijection.
\end{proof}
Together with CI and RL, this is a sufficient as well as necessary
test for the actual unit-\(\zeta\) first-square/second-cube branch.
A triple passing the test reconstructs the seed through \(t\).
No such triple is asserted to exist. The two signs of \(Z\)
retain the ordered seed directions.

\begin{theorem}[Integral factorization and exact prime depth]\label{ssg-integral}
Write a point of \eqref{ssg-positive-test} in primitive coordinates
\[
 X=A/d_0^2,\quad Y=B/d_0^3,\quad d_0>0,\quad
 \gcd(A,d_0)=\gcd(B,d_0)=1.
\]
Then \(C=Zd_0^3\in\mathbb Z\). With \(K=A+2d_0^2\ne0\),
\begin{align}
 B^2&=A^3-9Ad_0^4-9d_0^6,\notag\\
 (B^2-C^2)^2&=36d_0^2K^2(3B^2+C^2),\label{ssg-integral-conic}\\
 L:=\frac{B^2-C^2}{6d_0K}&\in\mathbb Z,\qquad
 L^2=3B^2+C^2.\notag
\end{align}
Every prime \(p>3\) dividing \(d_0K\) satisfies
\begin{equation}\label{ssg-exact-depth}
 v_p(B^2-C^2)=v_p(d_0K),\qquad p\nmid BCL.
\end{equation}
The entire \(p\)-primary part of \(d_0K\) divides exactly one of
\(B-C\) and \(B+C\), with no excess depth there.
\end{theorem}
\begin{proof}
The coordinates exist: at a prime of negative \(X\)-valuation,
the cubic term is uniquely smallest in the equation, so that
valuation is even and the \(Y\)-denominator exponent is three
times the corresponding half-exponent. Elsewhere \(Y\) is integral.
Both \(H_0=d_0^6H(A/d_0^2)\) and \(R_0=d_0^6R(A/d_0^2)\) are integers.
Clearing the quartic makes \(C\) a rational root of the monic polynomial
\[
 C^4-2H_0C^2+B^2R_0=0,
\]
so it is an integer. Clearing \eqref{ssg-conic} gives
\eqref{ssg-integral-conic}. The square of the nonzero integer \(6d_0K\)
divides \((B^2-C^2)^2\); prime valuations imply
\(6d_0K\mid B^2-C^2\), and hence the assertions about \(L\).

One has \(\gcd(d_0,K)=1\). If \(p\mid d_0\), then
\(B^2\equiv A^3\not\equiv0\pmod p\).
If \(p\mid K\), then \(d_0\) is a unit and
\(A\equiv-2d_0^2\), giving \(B^2\equiv d_0^6\not\equiv0\pmod p\).
In either case \eqref{ssg-integral-conic} forces
\(C^2\equiv B^2\pmod p\). Thus \(3B^2+C^2\equiv4B^2\) is a unit
for \(p>3\). Taking valuations proves \eqref{ssg-exact-depth},
and the equation for \(L\) gives \(p\nmid L\).
Such an odd prime cannot divide both \(B-C\) and \(B+C\);
their product is \(B^2-C^2\), proving the allocation.
\end{proof}
This necessary factorization gives no upper bound for \(A,d_0\)
or the point height, and does not force a prime violating the allocation.

\begin{theorem}[A three-isogeny between the elliptic quotients]\label{ssg-isogeny}
For the curves in \eqref{eqg-elliptic-curves}, put
\[
 F_3(X)=X+\frac{36}{X+3}-\frac{36}{(X+3)^2},\qquad
 F'_3(X)=1-\frac{36}{(X+3)^2}+\frac{72}{(X+3)^3}.
\]
The chart map
\begin{equation}\label{ssg-isogeny-map}
 \phi:E\longrightarrow E',\qquad
 (X,Y)\longmapsto(F_3(X),YF'_3(X))
\end{equation}
extends to a degree-three isogeny over \(\mathbb Q\), with geometric
kernel \(\{O,(-3,3i),(-3,-3i)\}\). Consequently
\(\operatorname{Jac}(H_1)\sim_{\mathbb Q}E\times E\).
\end{theorem}
\begin{proof}
Clearing \((X+3)^6\) verifies
\begin{equation}\label{ssg-isogeny-identity}
 f(X)F'_3(X)^2=F_3(X)^3-189F_3(X)+999.
\end{equation}
The nonconstant rational map extends between the smooth projective
curves. At \(O\), \(F_3(X)\) has leading term \(X\), so \(O\)
maps to \(O\). The origin-preserving morphism theorem
(Milne, \emph{Elliptic Curves}, Chapter II, Proposition 1.5)
makes it a homomorphism.
The numerator of \(F_3\) is cubic, its denominator is \((X+3)^2\),
and the numerator at \(-3\) is \(-36\). Thus its degree on
\(\mathbb P^1\) is three. Both elliptic abscissa maps have degree
two; their commutative diagram gives \(\deg\phi=3\).
The poles are exactly \(O\) and the two points \(X=-3,Y^2=-9\);
these are the preimages of \(O\).
Composition with the established EQ isogeny proves the final assertion.
\end{proof}
For \(P=(-2,1)\), \(P'=(6,9)\), one has
\[
 \phi(P)=(-2,37)=-3P'.
\]
Indeed \(2P'=(33/4,9/8)\); its secant with \(P'\) has slope
\(-7/2\), giving \(3P'=(-2,-37)\).
Neither this identity nor the isogeny certifies a rank upper bound,
a full Mordell--Weil basis, or the rational points passing
\eqref{ssg-positive-test}.

\paragraph{Verification and remaining scope.}
Both peer routes and the root group reviewed the ordinary proofs
in full. An exact standard-library replay checks fifteen polynomial
identities and the displayed point relation. It also checks the
eighty specified points \(nP\), \(1\leq |n|\leq40\), and their isogeny
images; none passes the positive gate. This bounded check is not a
complete group or rational-point computation.
The twelve declarations in \texttt{SimultaneousEllipticArithmetic}
also passed fresh compilation and all axiom queries with only the
standard axioms. They prove the actual polynomial identities,
the homogeneous elliptic/quartic-to-conic implication, integer
square divisibility and its quotient, and the boundary unit from
\(\operatorname{Int.gcd}(A,d_0)=1\). These are combined to prove
the full product-depth equality in \eqref{ssg-exact-depth} from
the actual homogeneous equations and stated prime hypotheses,
without a supplied depth or unit conclusion. The rational
three-isogeny identity is formalized as well.
The monic rational-to-integer bridge, primitive rational-coordinate
representation, projective point bijection, geometric isogeny
degree/kernel, Jacobian statement and final allocation to one of
the two linear factors remain ordinary-only.
The simultaneous rational-point problem remains unresolved, as do
the other nonidentity unit and the general ABC statement.
